\chapter{Subtraction}\label{ch:g2:subtraction}

When the units on top are too few to take from, we break open a
ten --- the \emph{borrow}, the mirror of the carry. This chapter
explains it with sticks, then trains the written subtraction, with
counting up as a friendly shortcut.

\section{Why we borrow}

\begin{example}[Unbundling a ten]\label{ex:g2:subtraction:unbundle}
$42 - 17$: we want to take $7$ units from $2$ units --- too few!
So untie one of the $4$ tens into ten loose units: now the top has
$3$ tens and $12$ units. Take away: $12 - 7 = 5$ units, and
$3 - 1 = 2$ tens. So $42 - 17 = 25$.
\end{example}

\begin{omfigure}
\begin{tikzpicture}[scale=0.5]
  \foreach \i in {0,...,9} \draw[omExo, very thick]
    (0.16*\i,0) -- (0.16*\i,1.4);
  \draw[-Stealth, omThm, thick] (2.0,0.7) -- (3.4,0.7);
  \foreach \i in {0,...,9} \draw[omDef, very thick]
    (4.0+0.22*\i,0) -- (4.0+0.22*\i,1.4);
  \node[below, font=\small] at (0.72,-0.35) {one ten};
  \node[below, font=\small] at (5.0,-0.35) {ten loose units};
  \node[right, font=\small] at (7.2,0.7)
    {untie a ten to have enough units to take away};
\end{tikzpicture}

{\small The borrow is unbundling: one ten becomes ten units, so the
subtraction can go through.}
\end{omfigure}

\section{The written method}

\begin{method}[Column subtraction]\label{met:g2:subtraction:column}
\begin{enumerate}
  \item Write the bigger number on top, aligned on the right;
  \item subtract the units; if the top digit is too small, borrow
        a ten (it becomes ten units in the current column);
  \item subtract the tens, then the hundreds;
  \item \emph{check}: the answer plus the number taken away must
        give back the top number.
\end{enumerate}
\end{method}

\begin{example}\label{ex:g2:subtraction:worked}
Compute $503 - 168$:
\[
\begin{array}{r}
5\,0\,3 \\
-\ 1\,6\,8 \\
\hline
3\,3\,5
\end{array}
\]
Units: $3 - 8$ too small; borrow through the tens (there are none,
so borrow from the hundreds first): the top becomes $4$ hundreds,
$9$ tens, $13$ units. Then $13 - 8 = 5$; $9 - 6 = 3$;
$4 - 1 = 3$. Check: $335 + 168 = 503$. \checkmark
\end{example}

\section{Counting up}

\begin{method}[Subtracting by counting up]\label{met:g2:subtraction:countup}
For numbers that are close, climb from the small one to the big
one and add the hops:
\[
83 - 78: \quad 78 \to 80 \text{ is } 2, \ 80 \to 83 \text{ is } 3,
\text{ so } 83 - 78 = 5 .
\]
Counting up is quicker than borrowing when the two numbers are near
each other --- and it is exactly how a shopkeeper gives change.
\end{method}

\begin{omfigure}
\begin{tikzpicture}[scale=0.85]
  \draw[-Stealth] (76.6,0) -- (84.4,0);
  \foreach \i in {77,...,84} \draw (\i,-0.1) -- (\i,0.1);
  \foreach \i in {78,80,83}
    \node[below, font=\footnotesize] at (\i,-0.12) {\i};
  \draw[-Stealth, omDef, thick] (78,0.3) .. controls (79,0.9) ..
    (80,0.3);
  \node[omDef, font=\small] at (79,1.0) {$+2$};
  \draw[-Stealth, omThm, thick] (80,0.3) .. controls (81.5,1.05) ..
    (83,0.3);
  \node[omThm, font=\small] at (81.5,1.1) {$+3$};
\end{tikzpicture}

{\small $83 - 78$ by counting up: two hops, $2 + 3 = 5$.}
\end{omfigure}

\begin{example}[Giving change]\label{ex:g2:subtraction:change}
A toy costs $34$; you pay with a $50$. Change: count up,
$34 \to 40$ is $6$, $40 \to 50$ is $10$, so the change is $16$.
The subtraction $50 - 34 = 16$, done the shopkeeper's way.
\end{example}

\section{Exercises}

\begin{exercise}[$\star$]\label{exo:g2:subtraction:1}
Compute (no borrow needed): $58 - 23$; \quad $176 - 45$; \quad
$389 - 164$.
\end{exercise}

\begin{exercise}[$\star$]\label{exo:g2:subtraction:2}
Compute with a borrow: $52 - 17$; \quad $73 - 28$; \quad
$60 - 24$.
\end{exercise}

\begin{exercise}[$\star$]\label{exo:g2:subtraction:3}
Compute in columns and check each one: $324 - 167$; \quad
$500 - 236$; \quad $703 - 158$.
\end{exercise}

\begin{exercise}[$\star$]\label{exo:g2:subtraction:4}
Compute by counting up: $62 - 58$; \quad $91 - 85$; \quad
$100 - 96$.
\end{exercise}

\begin{exercise}[$\star$]\label{exo:g2:subtraction:5}
Compute the change (count up): pay $50$ for a $42$ book; pay $20$
for a $13$ toy; pay $100$ for a $75$ bag.
\end{exercise}

\begin{exercise}[$\star$]\label{exo:g2:subtraction:6}
``A tree is $156$ cm tall. It grew $28$ cm this year. How tall was
it last year?'' Write the subtraction and answer.
\end{exercise}

\begin{exercise}[$\star$]\label{exo:g2:subtraction:7}
Complete with the twin addition to check: $85 - 39 = 46$; is it
right? $214 - 68 = 156$; is it right?
\end{exercise}

\begin{exercise}[$\star$]\label{exo:g2:subtraction:8}
For each story, choose $+$ or $-$: ``$47$ birds, $19$ fly away'';
``$47$ apples, $19$ more picked''; ``Sam has $47$, Lea has $19$:
how many more has Sam?''.
\end{exercise}

\begin{exercise}[$\star$]\label{exo:g2:subtraction:9}
Fill the hole: $63 - \;?\; = 45$. (Count up from $45$ to $63$.)
\end{exercise}

\begin{exercise}[$\star\star$]\label{exo:g2:subtraction:10}
A game: think of a number, add $30$, then take away $18$. The
result is $58$. What was the number? (Work backward!)
\end{exercise}
